% !TeX program = xelatex
% 配套原卷：同目录“专题 基本不等式.pdf”，5页、9类、33题。
% 本文件是独立推导的解析稿；原卷及其题号保持不变。
\documentclass[UTF8,a4paper,11pt]{ctexart}
\usepackage[margin=20mm,headheight=15pt]{geometry}
\usepackage{amsmath,amssymb,booktabs,array,fancyhdr,xcolor,needspace}
\usepackage{newtxtext,newtxmath}
\usepackage[hidelinks]{hyperref}
\definecolor{accent}{HTML}{24566E}
\pagestyle{fancy}\fancyhf{}
\fancyhead[L]{\small 基本不等式求最值}\fancyhead[R]{\small 九类经典题 · 逐题详解}
\fancyfoot[C]{\small\thepage}
\setlength{\parindent}{0pt}
\setlength{\parskip}{4pt}
\setlength{\abovedisplayskip}{6pt}
\setlength{\belowdisplayskip}{6pt}
\ctexset{section={format=\Large\bfseries\color{accent},beforeskip=0pt,afterskip=10pt}}
\newcommand{\topic}[2]{\clearpage\section*{#1}\vspace{-4pt}{\small\color{gray}原卷位置：第#2页。题号与原卷一致。}\par\medskip}
\newcommand{\q}[3]{\par\Needspace{7\baselineskip}\textbf{\color{accent}第#1题}\quad #2\par\textbf{答案：}#3\par}
\newcommand{\tip}[1]{\par{\small\textbf{易错提醒：}#1}\par}
\begin{document}
\begin{center}
{\LARGE\bfseries 专题：基本不等式}\par\vspace{8pt}
{\Large 九类经典题 · 33题逐题详解}\par\vspace{6pt}
{\small 配套原卷《专题 基本不等式.pdf》}
\end{center}
\vspace{5pt}
\textbf{使用说明}\quad 按原卷题型、题号逐一对应，保留全部33题。各题先列核心题设及答案，再给推导和合法取等条件；选择题字母对应原卷选项。本册为配套解析，不替代原始题目。

\textbf{来源说明}\quad 原卷位于本册同目录，共5页，页脚标注“学科网（北京）股份有限公司”。本册逐页核对题干并独立推导答案；原卷标注的年份、地区及考试名称未经外部核验，不作为已核实真题出处。

\section*{答案速查}
\renewcommand{\arraystretch}{1.7}
\begin{tabular}{@{}p{0.28\textwidth}p{0.67\textwidth}@{}}
\toprule
题型 & 按原题顺序列出的答案 \\
\midrule
一、直接法 & （1）D；（2）A；（3）D；（4）D；（5）$\dfrac1{16}$ \\
二、配凑法 & （1）C；（2）A；（3）D \\
三、分式 & （1）$2\sqrt2-3$；（2）$2\sqrt2$；（3）A；（4）D \\
四、常数代换 & （1）D；（2）D；（3）A；（4）C \\
五、换元法 & （1）C；（2）B；（3）B；（4）$31$ \\
六、消元法 & （1）D；（2）A；（3）C；（4）A \\
七、和积共存 & （1）$xy\in[16,+\infty)$，$x+y\in[8,+\infty)$；（2）D；（3）D \\
八、恒成立与有解 & （1）B；（2）A；（3）A \\
九、实际应用 & （1）C；（2）A；（3）D \\
\bottomrule
\end{tabular}\par\vspace{6pt}

\section*{每题都要检查的三件事}
\textbf{一正：}使用 $u+v\ge2\sqrt{uv}$ 前，确认 $u,v\ge0$；有倒数时分母不能为零。

\textbf{二定：}求和的下界时看积是否固定，求积的上界时看和是否固定。必要时先配凑、换元或消元。

\textbf{三相等：}等号要求本次比较的两个量相等，还要同时满足原题条件。不能把不能取到的下界写成最小值。

\textbf{特别关注：}题型二第3题要处理全体实数；题型四第2题、题型七第1题要证明完整值域；题型六第3题要检查换元范围；题型七第3题要区分开端点、最小值与最大值；题型八要区分“存在”与“任意”。

\topic{题型一：直接法求最值}{1}
\q{1}{$x>0$，求 $y=3x+\dfrac1x$ 的最小值。}{D，$2\sqrt3$。}
因为 $3x>0$、$1/x>0$，且二者乘积为3，所以
\[
3x+\frac1x\ge2\sqrt{3x\cdot\frac1x}=2\sqrt3.
\]
等号成立当且仅当 $3x=1/x$，即 $x=1/\sqrt3$，符合 $x>0$，故最小值为 $2\sqrt3$。

\q{2}{$x,y$ 为非零实数，求 $\dfrac{y^2}{x^2}+\dfrac{9x^2}{y^2}$ 的最小值。}{A，$6$。}
虽然 $x,y$ 不一定为正，但两个平方比值均为正，且乘积为9，因此
\[
\frac{y^2}{x^2}+\frac{9x^2}{y^2}\ge2\sqrt9=6.
\]
等号条件为 $y^4=9x^4$，即 $y^2=3x^2$，也就是 $y=\pm\sqrt3x$（$x\ne0$）。例如 $x=1,y=\sqrt3$ 时取到6。

\q{3}{$x>0$，求 $y=(1-x)\left(8-\dfrac2x\right)$ 的最大值。}{D，$2$。}
先展开，再把问题转化为正数和的最小值：
\[
y=10-\left(8x+\frac2x\right)\le10-2\sqrt{8x\cdot\frac2x}=2.
\]
等号成立当且仅当 $8x=2/x$，即 $x=1/2$。代回原式得 $y=(1/2)\times4=2$。
\tip{原式两个因子可能为负，不能直接把它们作为正数套用基本不等式。}

\q{4}{实数 $m,n$ 满足 $m^2+n^2=8$，求 $mn$ 的最大值。}{D，$4$。}
由 $(m-n)^2\ge0$ 得 $2mn\le m^2+n^2=8$，所以 $mn\le4$。
等号成立当且仅当 $m=n$，结合条件得 $(m,n)=(2,2)$ 或 $(-2,-2)$。因此最大值为4。
\tip{此处用平方非负最稳妥；原题没有要求 $m,n$ 为正数。}

\q{5}{正实数 $a,b$ 满足 $a+4b=1$，求 $ab$ 的最大值。}{$\dfrac1{16}$。}
把 $a$ 与 $4b$ 作为一对正数，有
\[
1=(a+4b)^2\ge4a\cdot4b=16ab.
\]
故 $ab\le1/16$。等号成立当且仅当 $a=4b$，联立 $a+4b=1$，得 $a=1/2,b=1/8$，确能取到。

\topic{题型二：配凑法求最值}{1}
\q{1}{$a>1$，求 $4a+\dfrac1{a-1}$ 的最小值。}{C，$8$。}
分母是 $a-1$，因此将 $4a$ 拆成 $4(a-1)+4$。由于 $a-1>0$，
\[
4a+\frac1{a-1}=4+4(a-1)+\frac1{a-1}
\ge4+2\sqrt4=8.
\]
等号成立当且仅当 $4(a-1)=1/(a-1)$，由 $a>1$ 得 $a-1=1/2$，即 $a=3/2$。代入原式得 $6+2=8$。
\tip{直接比较 $4a$ 与 $1/(a-1)$，两者乘积不是定值，通常不能得到所求最小值。}

\q{2}{$x>0$，求 $y=x+\dfrac12+\dfrac1{2x+1}$ 的最小值。}{A，$\sqrt2$。}
把前两项合成分母的倍数：
\[
y=\frac{2x+1}{2}+\frac1{2x+1}
\ge2\sqrt{\frac{2x+1}{2}\cdot\frac1{2x+1}}=\sqrt2.
\]
等号条件为 $(2x+1)^2=2$。由 $2x+1>1$ 得
\[
2x+1=\sqrt2,\qquad x=\frac{\sqrt2-1}{2}>0.
\]
所以最小值确为 $\sqrt2$。
\tip{若令 $t=2x+1$，其范围是 $t>1$，不只是 $t>0$；取等点 $t=\sqrt2$ 位于该范围内。}

\q{3}{求函数 $y=x(3-2x)$ 的最大值（原题未另限定义域）。}{D，$\dfrac98$。}
\textbf{方法一：配方，直接覆盖所有实数。}
\[
y=-2x^2+3x=-2\left(x-\frac34\right)^2+\frac98\le\frac98.
\]
当且仅当 $x=3/4$ 时取等号，因此最大值为 $9/8$。

\textbf{方法二：使用基本不等式时先分范围。}
当 $0\le x\le3/2$ 时，$2x\ge0$、$3-2x\ge0$，且两者之和为3，所以
\[
2y=(2x)(3-2x)\le\left(\frac{2x+3-2x}{2}\right)^2=\frac94.
\]
等号成立当且仅当 $2x=3-2x$，即 $x=3/4$。
当 $x<0$ 或 $x>3/2$ 时，$x$ 与 $3-2x$ 异号，故 $y<0<9/8$。合并可知全局最大值为 $9/8$。
\tip{不能在全体实数上直接断言 $2x$ 与 $3-2x$ 都为正。}

\topic{题型三：分式求最值}{2}
\q{1}{$x>0$，求 $\dfrac{2x^2-3x+1}{x}$ 的最小值。}{$2\sqrt2-3$。}
分项相除，使定积结构显现：
\[
\frac{2x^2-3x+1}{x}=2x+\frac1x-3\ge2\sqrt2-3.
\]
等号成立当且仅当 $2x=1/x$，即 $x=1/\sqrt2>0$。
\tip{基本不等式作用于 $2x$ 与 $1/x$；原分式最小值为负并不矛盾。}

\q{2}{$x>1$，求 $\dfrac{x^2-2x+3}{x-1}$ 的最小值。}{$2\sqrt2$。}
注意 $x^2-2x+3=(x-1)^2+2$，于是
\[
\frac{x^2-2x+3}{x-1}=(x-1)+\frac2{x-1}\ge2\sqrt2.
\]
这里 $x-1>0$，等号成立当且仅当 $(x-1)^2=2$，即 $x=1+\sqrt2$，符合原条件。

\q{3}{$x<-1$，求 $\dfrac{x^2-x+14}{x+1}$ 的最大值。}{A，$-11$。}
先作除法，再将负量换成正量：
\[
\frac{x^2-x+14}{x+1}=x-2+\frac{16}{x+1}.
\]
令 $t=-(x+1)>0$，则 $x-2=-t-3$，所以
\[
\frac{x^2-x+14}{x+1}=-3-\left(t+\frac{16}{t}\right)
\le-3-8=-11.
\]
等号成立当且仅当 $t=16/t$，即 $t=4$，对应 $x=-5<-1$。代回原式为 $(25+5+14)/(-4)=-11$，因此最大值为 $-11$。
\tip{应先处理 $x+1<0$ 的符号。负的“正数和”取得最大值，恰在该正数和取得最小值时。}

\q{4}{正实数 $x,y,z$ 满足 $x^2-xy+y^2-z=0$，求 $xy/z$ 的最大值。}{D，$1$。}
由条件可得
\[
z=x^2-xy+y^2=(x-y)^2+xy\ge xy>0.
\]
分子、分母均为正，因此 $xy/z\le1$。当且仅当 $x=y>0$ 时等号成立，此时 $z=x^2$。例如 $x=y=z=1$ 满足条件，故最大值为1。

\topic{题型四：常数代换法求最值}{2--3}
\q{1}{$x,y>0$，$x+y=1$，求 $\dfrac1x+\dfrac4y$ 的最小值。}{D，$9$。}
用 $x+y$ 代换常数1，再展开：
\[
\frac1x+\frac4y=\left(\frac1x+\frac4y\right)(x+y)
=5+\frac yx+\frac{4x}y\ge5+4=9.
\]
等号条件为 $y/x=4x/y$。因 $x,y>0$，即 $y=2x$；结合 $x+y=1$，得 $x=1/3,y=2/3$。

\q{2}{$m,n>0$，$2m+n=1$，求 $\dfrac nm+\dfrac1n$ 的取值范围。}{D，$[1+2\sqrt2,+\infty)$。}
因为 $1/n=(2m+n)/n=1+2m/n$，所以
\[
\frac nm+\frac1n=1+\frac nm+\frac{2m}n\ge1+2\sqrt2.
\]
等号成立当且仅当 $n^2=2m^2$，即 $n=\sqrt2m$。联立得
$m=1/(2+\sqrt2)$、$n=\sqrt2-1$，均为正。

\textbf{证明整个区间都能取到：}令 $t=n/m>0$，则 $m=1/(2+t)$、$n=t/(2+t)$，任何 $t>0$ 都给出合法的 $m,n$。目标式成为 $1+t+2/t$。
任给 $k\ge1+2\sqrt2$，方程 $t^2-(k-1)t+2=0$ 的判别式非负，两根的和为正、积为正，故有正根。于是每个这样的 $k$ 均能取到。

\q{3}{$x,y>0$，$x+2y-2xy=0$，求 $x+2y$ 的最小值。}{A，$4$。}
原条件等价于 $x+2y=2xy$。设 $s=x+2y>0$，则
\[
s=x+2y\ge2\sqrt{2xy}=2\sqrt s.
\]
两边除以 $\sqrt s>0$ 得 $\sqrt s\ge2$，即 $s\ge4$。
等号成立要求 $x=2y$，代入原约束得 $y=1,x=2$，确有 $x+2y=4$。
\tip{约束中同时出现和与积时，可先把乘积换成目标式，再解关于目标式的不等式。}

\q{4}{$a,b>0$，$a+\dfrac1b=1$，求 $\dfrac1a+b$ 的最小值。}{C，$4$。}
利用原条件乘以1：
\[
\frac1a+b=\left(\frac1a+b\right)\left(a+\frac1b\right)
=2+ab+\frac1{ab}\ge4.
\]
等号成立当且仅当 $ab=1$。此时 $1/b=a$，由 $a+1/b=1$ 得 $a=1/2,b=2$，均符合原条件。

\topic{题型五：换元法求最值}{3}
\q{1}{$a,b>0$，$a+b=2$，求 $\dfrac4{1+a}+\dfrac1{1+b}$ 的最小值。}{C，$\dfrac94$。}
令 $u=1+a,v=1+b$，则 $u,v>1$，$u+v=4$。因此
\[
\frac4u+\frac1v=\frac14\left(\frac4u+\frac1v\right)(u+v)
=\frac14\left(5+\frac{4v}u+\frac uv\right)\ge\frac94.
\]
等号成立当且仅当 $u=2v$。结合 $u+v=4$ 得 $u=8/3,v=4/3$，对应 $a=5/3,b=1/3$，满足正数条件。

\q{2}{$0<x<1$，求 $\dfrac1x+\dfrac{16}{1-x}$ 的最小值。}{B，$25$。}
令 $u=x,v=1-x$，则 $u,v>0$、$u+v=1$，从而
\[
\frac1u+\frac{16}v=\left(\frac1u+\frac{16}v\right)(u+v)
=17+\frac vu+\frac{16u}v\ge25.
\]
等号成立当且仅当 $v=4u$，即 $1-x=4x$，所以 $x=1/5\in(0,1)$。

\q{3}{$a,b>0$，$2a+b=3$，求 $\dfrac{4a^2+1}{2a}+\dfrac{b^2+2b+2}{b+1}$ 的最小值。}{B，$5$。}
先拆分，注意 $b^2+2b+2=(b+1)^2+1$：
\[
\frac{4a^2+1}{2a}+\frac{b^2+2b+2}{b+1}
=2a+(b+1)+\frac1{2a}+\frac1{b+1}.
\]
令 $u=2a>0,v=b+1>1$，则 $u+v=4$，原式为 $4+1/u+1/v$。
由 $uv\le(u+v)^2/4=4$，得
\[
4+\frac1u+\frac1v=4+\frac4{uv}\ge5.
\]
等号成立当且仅当 $u=v=2$，即 $a=b=1$，满足原条件。

\q{4}{$x,y>0$，$\dfrac1{x+1}+\dfrac3{x+y}=\dfrac12$，求 $4x+3y$ 的最小值。}{$31$。}
令 $u=x+1,v=x+y$，则 $u,v>0$，$1/u+3/v=1/2$，且 $4x+3y=u+3v-1$。于是
\[
\frac{u+3v}{2}=(u+3v)\left(\frac1u+\frac3v\right)
=10+3\left(\frac uv+\frac vu\right)\ge16.
\]
故 $u+3v\ge32$，$4x+3y\ge31$。等号成立时 $u=v$，由 $4/u=1/2$ 得 $u=v=8$，还原为 $x=7,y=1>0$，合法且取到31。

\topic{题型六：消元法求最值}{3--4}
\q{1}{$x,y>0$，$x+\dfrac1y=1$，求 $\dfrac1x+2y$ 的最小值。}{D，$3+2\sqrt2$。}
由条件得 $x=1-1/y>0$，因此 $y>1$。消去 $x$：
\[
\frac1x+2y=\frac y{y-1}+2y
=3+2(y-1)+\frac1{y-1}\ge3+2\sqrt2.
\]
等号条件为 $2(y-1)^2=1$，故 $y=1+1/\sqrt2$，$x=\sqrt2-1>0$，满足原约束。

\q{2}{$x,y>0$，$x^2+3xy-2=0$，求 $2x+y$ 的最小值。}{A，$\dfrac{2\sqrt{10}}3$。}
由 $3xy=2-x^2>0$ 得 $0<x<\sqrt2$，并且 $y=(2-x^2)/(3x)$，所以
\[
2x+y=\frac{5x}{3}+\frac2{3x}\ge2\sqrt{\frac{5x}{3}\cdot\frac2{3x}}=\frac{2\sqrt{10}}3.
\]
等号成立当且仅当 $5x^2=2$，即 $x=\sqrt{2/5}=\sqrt{10}/5$，位于合法范围内；此时 $y=4\sqrt{10}/15>0$，代回约束成立。

\q{3}{$m,n>0$，$mn=2$，求 $\dfrac1m+\dfrac2n+\dfrac9{2m+n}$ 的最小值。}{C，$3\sqrt2$。}
由 $mn=2$ 得 $1/m=n/2$、$2/n=m$，故前两项之和为 $(2m+n)/2$。
令 $t=2m+n$，则
\[
t\ge2\sqrt{2mn}=4,\qquad
\frac1m+\frac2n+\frac9{2m+n}=\frac t2+\frac9t\ge3\sqrt2.
\]
等号成立当且仅当 $t^2=18$，即 $t=3\sqrt2>4$。还须确认原变量能实现它：联立 $2m+n=3\sqrt2$、$mn=2$，可得
\[
(m,n)=(\sqrt2,\sqrt2)\quad\text{或}\quad
\left(\frac1{\sqrt2},2\sqrt2\right).
\]
两组都为正且满足原条件，所以最小值为 $3\sqrt2$。
\tip{$t$ 的最小值为4，并不表示 $t/2+9/t$ 在 $t=4$ 取得最小值。}

\q{4}{$b>0$，$ab+b=1$，求 $a^2b$ 的最小值。}{A，$0$。}
因为 $a^2\ge0$、$b>0$，直接有 $a^2b\ge0$。取 $a=0,b=1$ 即满足原约束并取到0。
也可消元：$a=1/b-1$，故 $a^2b=b+1/b-2\ge0$，等号在 $b=1,a=0$ 时成立。
\tip{题设仅要求 $b>0$，没有要求 $a>0$，不能自行排除 $a=0$。}

\topic{题型七：和积共存求最值}{4}
\q{1}{正数 $x,y$ 满足 $x+y+8=xy$，求 $xy$ 与 $x+y$ 的取值范围。}{$xy\in[16,+\infty)$，$x+y\in[8,+\infty)$。}
\textbf{（1）求 $xy$ 的范围。}设 $p=xy>0$，则 $x+y=p-8$，由基本不等式
\[
p-8=x+y\ge2\sqrt p.
\]
令 $t=\sqrt p>0$，则 $t^2-2t-8=(t-4)(t+2)\ge0$。因 $t+2>0$，有 $t\ge4$，所以 $p\ge16$。等号在 $x=y=4$ 时成立。

\textbf{（2）求 $x+y$ 的范围。}由 $x+y=p-8$，得到 $x+y\ge8$。

\textbf{充分性：}任给 $s\ge8$，令 $x,y$ 为方程 $z^2-sz+(s+8)=0$ 的两根。
其判别式为 $\Delta=s^2-4s-32=(s-8)(s+4)\ge0$，两根的和 $s>0$、积 $s+8>0$，故两根均为正，并满足 $x+y+8=xy$。因此每个 $s\ge8$ 都可取到，相应的每个 $p=s+8\ge16$ 也都可取到。

\q{2}{$a\ge b>0$，$ab=a+b+15$，求 $ab$ 的最小值。}{D，$25$。}
令 $t=\sqrt{ab}>0$，由 $a+b\ge2\sqrt{ab}$ 得
\[
t^2-15\ge2t\ \Longrightarrow\ (t-5)(t+3)\ge0.
\]
所以 $t\ge5$，$ab\ge25$。等号成立当且仅当 $a=b$，代回得 $a=b=5$，符合 $a\ge b>0$，因此最小值为25。

\q{3}{$a,b>0$、$b>\dfrac12$，且 $a+2b=2ab-3$，判断四个选项。}{D，$2a+b$ 的最小值为 $\dfrac{13}2$。}
整理条件得
\[
(a-1)(2b-1)=4.
\]
因 $2b-1>0$，必有 $a>1$。令 $u=a-1>0,v=2b-1>0$，则 $uv=4$。

\textbf{A错误：}任意 $u>0$ 都可令 $v=4/u>0$，故 $a$ 的范围是 $(1,+\infty)$，不含1。

\textbf{B错误：}$a+1/a-2=(a-1)^2/a>0$，所以不能取到2；当 $a$ 趋近于1的右侧时可任意接近2，因此2只是下确界，没有最小值。

\textbf{C错误：}$ab=(u+1)(v+1)/2=(5+u+v)/2\ge9/2$，等号在 $u=v=2$，即 $a=3,b=3/2$ 时成立。$9/2$ 是最小值；令 $u$ 增大、$v=4/u$，可知 $ab$ 无上界。

\textbf{D正确：}
\[
2a+b=\frac52+2u+\frac v2\ge\frac52+2\sqrt{uv}=\frac{13}2.
\]
等号要求 $2u=v/2$，结合 $uv=4$ 得 $u=1,v=4$，即 $a=2,b=5/2$，满足原题所有条件。

\topic{题型八：基本不等式的恒成立、有解问题}{5}
\q{1}{正实数 $x,y$ 满足 $\dfrac1x+\dfrac4y=1$，若 $x+\dfrac y4<m$ 有解，求 $m$ 的范围。}{B，$m>4$。}
令 $u=x>0,v=y/4>0$，则 $1/u+1/v=1$，目标式为 $u+v$。
由
\[
u+v=(u+v)\left(\frac1u+\frac1v\right)
=2+\frac uv+\frac vu\ge4,
\]
等号成立当且仅当 $u=v$，结合约束得 $u=v=2$，即 $x=2,y=8$。

若 $m\le4$，则总有 $x+y/4\ge4\ge m$，严格不等式无解；若 $m>4$，取 $x=2,y=8$ 就有 $x+y/4=4<m$。所以所求范围恰为 $(4,+\infty)$。
\tip{“有解”要求存在一组合法变量；本题是严格小于，故 $m=4$ 必须排除。}

\q{2}{$x,y>0$、$x+y=5$，若 $\dfrac4{x+1}+\dfrac1{y+2}\ge2m+1$ 恒成立，求 $m$ 的范围。}{A，$m\le\dfrac1{16}$。}
令 $u=x+1>1,v=y+2>2$，则 $u+v=8$。由
\[
\frac4u+\frac1v=\frac18\left(\frac4u+\frac1v\right)(u+v)
=\frac18\left(5+\frac{4v}u+\frac uv\right)\ge\frac98,
\]
等号要求 $u=2v$。由 $u+v=8$ 得 $u=16/3,v=8/3$，即 $x=13/3,y=2/3$，均为正，所以左端的最小值为 $9/8$。

原不等式对所有合法的 $x,y$ 恒成立，当且仅当
\[
2m+1\le\frac98\ \Longleftrightarrow\ m\le\frac1{16}.
\]
当 $m=1/16$ 时仍允许等号，所以端点包含在内。

\q{3}{$x,y>0$、$x+2y=1$，若 $\dfrac2x+\dfrac1y\ge m^2+7m$ 恒成立，求 $m$ 的范围。}{A，$-8\le m\le1$。}
用约束代换1并展开：
\[
\frac2x+\frac1y=\left(\frac2x+\frac1y\right)(x+2y)
=4+\frac{4y}x+\frac xy\ge8.
\]
等号要求 $4y/x=x/y$，即 $x=2y$，结合 $x+2y=1$ 得 $x=1/2,y=1/4$，所以最小值为8。
因此恒成立等价于
\[
m^2+7m\le8\ \Longleftrightarrow\ (m+8)(m-1)\le0,
\]
解得 $m\in[-8,1]$，两个端点均可取。

\topic{题型九：利用基本不等式解决实际问题}{5}
\q{1}{用两臂不等长的天平，将5g砝码分别放在左右盘，各称一次黄金交付，判断总质量。}{C，黄金多给了。}
设左、右臂长分别为 $l,r>0$，且 $l\ne r$。设第一次称得黄金质量为 $m_1$，第二次为 $m_2$（单位均为g）。由天平平衡时两边力矩相等，约去相同的重力加速度后，得到
\[
5l=m_1r,\qquad m_2l=5r.
\]
故总质量为
\[
M=m_1+m_2=5\left(\frac lr+\frac rl\right)\ge10.
\]
等号成立当且仅当 $l=r$，但题设明确两臂不等长，所以实际为 $M>10$。两次称量的先后顺序不改变总质量。
\tip{“两次误差抵消”并不成立。这里得到的是严格大于10g，关键是等号条件被题设排除了。}

\q{2}{昆虫跳跃高度 $H$ 与水平速度 $v$ 满足 $v^2=\dfrac{4H}{1-Hv^2}$，求最大跳跃高度。}{A，$0.25$米。}
高度非负，且原式要求 $1-Hv^2\ne0$。交叉相乘并整理：
\[
v^2-Hv^4=4H\ \Longrightarrow\ H=\frac{v^2}{v^4+4}.
\]
当 $v=0$ 时，$H=0$。当 $v\ne0$ 时，$v^2>0$，可写成
\[
H=\frac1{v^2+4/v^2}\le\frac1{2\sqrt4}=\frac14.
\]
等号成立当且仅当 $v^2=4/v^2$，即 $v^2=2$。按速度为非负数，取 $v=\sqrt2$ 米/秒，此时 $H=1/4$ 米。
还需检查原分母：$1-Hv^2=1-1/2=1/2\ne0$，所以取等点合法。最大高度为 $0.25$ 米。

\q{3}{17列货车以 $v$ km/h 匀速从A到B，路程400km，相邻车距至少 $(v/20)^2$ km，求全部运到的最短时间。}{D，$8$小时。}
从首列货车由A发车开始计时，到第17列到达B结束。忽略车长、加速和调度耗时，最快时各相邻车距取允许的最小值。
17列车共有16个间隔，每个间隔对应的发车时间差为
\[
\frac{(v/20)^2}{v}=\frac v{400}\ \text{小时}\qquad(v>0).
\]
末列比首列晚发车 $16v/400=v/25$ 小时，末列本身还需行驶 $400/v$ 小时，所以总时间
\[
T=\frac{400}{v}+\frac v{25}\ge2\sqrt{\frac{400}{25}}=8\ \text{小时}.
\]
等号要求 $400/v=v/25$，得 $v=100$ km/h。此时相邻车距25km、发车间隔 $1/4$ 小时；首末车发车相差4小时，行程用时4小时，共8小时。
\tip{17列车之间只有16个间隔；只算单列运行时间 $400/v$ 会漏掉末列的发车等待。}
\end{document}
